
\section{弗洛尔上同调回顾 I：无反常}
\label{sec:2-3}

本节将对每个给定的横截对子 $(L_0,L_1)$，研究作用量泛函
$\mathcal A:\widetilde\Omega(L_0,L_1;\ell_0)\to\mathbb R$ 的有界梯度流线
$u:\mathbb R\to\widetilde\Omega(L_0,L_1;\ell_0)$ 所成的集合。由于定义域
$\mathbb R\times[0,1]$ 非紧，为了研究解 $u$ 的紧性性质与形变理论，需要对
$u$ 的导数施加某种衰减条件。我们将通过要求流线满足某种有界性来做到这一点；
流线的有界性由 $u$ 的辛面积来度量。
\translatornote{原书此处把 $\mathcal A$ 及梯度流线放在普通道路空间 $\Omega$ 上；2.2 节只在诺维科夫覆盖 $\widetilde\Omega$ 上定义了该泛函，译文依其类型补正。}

\subsection{作用量泛函的 \texorpdfstring{$L^2$}{L2}-梯度方程}
\label{sec:2-3-1}

回忆恒等式
\[
  d\mathcal A=-\pi^*\alpha.
\]
若通过覆盖投影把切空间
$T_{[\ell,w]}\widetilde\Omega(L_0,L_1;\ell_0)$ 与
$T_\ell\Omega(L_0,L_1)$ 等同起来，就可以在诺维科夫覆盖空间
$\widetilde\Omega(L_0,L_1;\ell_0)$ 上把这个恒等式写成
\begin{equation}
  d\mathcal A([\ell,w])(\xi)
  =\int_0^1\omega\bigl(\xi(t),\dot\ell(t)\bigr)\,dt.
  \tag{2.3.1}\label{eq:2-3-1}
\end{equation}
与有限维莫尔斯理论一样，我们将借助
$\widetilde\Omega(L_0,L_1;\ell_0)$ 上给定的一个“黎曼度量”来研究
$\mathcal A$ 的梯度流。

对每个给定的对子 $(L_0,L_1)$ 以及 $(M,\omega)$ 上一对相容殆复结构
$J_0,J_1$，考虑连接二者的 $(M,\omega)$ 上相容殆复结构道路
$\{J_t\}_{0\leq t\leq1}$。记
\[
  \mathcal J_\omega=\mathcal J_\omega(M)
  =\text{相容殆复结构所成的空间},
\]
以及
\[
  \begin{split}
  \mathcal P(\mathcal J_\omega)
  &=C^\infty([0,1],\mathcal J_\omega)\\
  &=\text{带参数相容殆复结构所成的空间}.
  \end{split}
\]
用 $\{J_t\}_t$ 表示 $\mathcal P(\mathcal J_\omega)$ 的一个元素，并令
\translatornote{原书把道路端点条件排成恒等式 $J_0=J_0,\ J_1=J_1$；这里展开参数位置，以区分道路取值与预先固定的端点。}
\begin{equation}
  \mathcal P(\mathcal J_\omega;J_0,J_1)
  =\bigl\{\{J_t\}_t\in\mathcal P(\mathcal J_\omega)
    \bigm|J_{t=0}=J_0,\ J_{t=1}=J_1\bigr\}.
  \tag{2.3.2}\label{eq:2-3-2}
\end{equation}
于是，对每个给定的 $\{J_t\}_t$，利用 $M$ 上的单参数黎曼度量族
\[
  g_t:=\omega(\mathord\cdot,J_t\mathord\cdot),
\]
由下式在 $\Omega$ 上定义一个 $L^2$ 度量：
\[
  \langle\xi_1,\xi_2\rangle_{\{J_t\}_t}
  :=\int_0^1\omega\bigl(\xi_1(t),J_t\xi_2(t)\bigr)\,dt.
\]
由式~\eqref{eq:2-3-1}，作用量泛函 $\mathcal A$ 的 $L^2$-梯度方程\liftIndex{I041}变为
\begin{equation}
  \begin{cases}
    \displaystyle\frac{\partial u}{\partial\tau}
      +J_t\frac{\partial u}{\partial t}=0,\\[6pt]
    u(\tau,0)\in L_0,
      \qquad u(\tau,1)\in L_1,
  \end{cases}
  \tag{2.3.3}\label{eq:2-3-3}
\end{equation}
其中 $u:\mathbb R\times[0,1]\to M$，并把 $u$ 看成
$\Omega(L_0,L_1;\ell_0)$ 中的道路 $\tau\mapsto u(\tau)$；这里道路
$u(\tau)$ 由 $u(\tau)(t)=u(\tau,t)$ 定义。

从分析观点看，要研究式~\eqref{eq:2-3-3} 之解的形变理论，控制 $u$ 的几何能量
\[
  E(u;\{J_t\}_t)
  :=\frac12\int_{\mathbb R}\int_0^1
    \left(
      \left|\frac{\partial u}{\partial\tau}\right|_{J_t}^{2}
      +\left|\frac{\partial u}{\partial t}\right|_{J_t}^{2}
    \right)dt\,d\tau
\]
至关重要。另一方面，只要 $\{J_t\}_t$ 是相容殆复结构族，对满足
\[
  \frac{\partial u}{\partial\tau}
  +J_t\frac{\partial u}{\partial t}=0
\]
的映射 $u$，便有恒等式
\begin{equation}
  E(u;\{J_t\}_t)=\int u^*\omega.
  \tag{2.3.4}\label{eq:2-3-4}
\end{equation}
用物理学语言来说，式~\eqref{eq:2-3-4} 在壳上成立。注意，辛面积
$\int u^*\omega$ 在同伦下不变。

现在设 $L_0,L_1$ 是 $(M,\omega)$ 中相互横截的拉格朗日子流形。对每个给定的
$\{J_t\}_t\in\mathcal P(\mathcal J_\omega)$，研究柯西--黎曼方程~\eqref{eq:2-3-3}，并把它的有界解空间定义为
\begin{equation}
  \widetilde{\mathcal M}^{\mathrm{reg}}
    (L_1,L_0;\{J_t\}_t)
  =\left\{
    u\ \middle|\
    \begin{array}{l}
      u\text{ 满足式~\eqref{eq:2-3-3}，并且}\\[-2pt]
      \displaystyle\int_{\mathbb R\times[0,1]}u^*\omega<\infty
    \end{array}
  \right\}.
  \tag{2.3.5}\label{eq:2-3-5}
\end{equation}
这里，记号中的 ``reg'' 代表 ``regular''（正则）。采用这个记号，是因为该模空间由
从 $\mathbb R\times[0,1]$ 出发的伪全纯映射组成。

有时我们会从各种记号中省略 $\{J_t\}_t$。还要说明，在模空间的记号中，我们把
$L_1$ 放在 $L_0$ 之前；这纯粹是我们的记号选择。

下面的命题给出自动的指数收敛结论：只要 $L_0$ 与 $L_1$ 横截，对任意
$u\in\widetilde{\mathcal M}^{\mathrm{reg}}(L_1,L_0)$，当
$\tau\to\pm\infty$ 时都有这种收敛。证明见\textbf{[Flo88IV]}。

\begin{proposition}{2.3.6}\label{prop:2-3-6}
设 $L_0,L_1$ 是横截相交的紧拉格朗日子流形。则对式~\eqref{eq:2-3-3}
任意辛面积有限的解 $u$，有极限
\[
  \lim_{\tau\to\infty}u(-\tau,\mathord\cdot)=\ell_p,
  \qquad
  \lim_{\tau\to\infty}u(\tau,\mathord\cdot)=\ell_q,
\]
其中 $\ell_p,\ell_q$ 分别是对应于某些唯一的
$p,q\in L_0\cap L_1$ 的常值道路。
\end{proposition}

这个命题给出分解
\[
  \widetilde{\mathcal M}^{\mathrm{reg}}(L_1,L_0)
  =\bigcup_{p,q\in L_0\cap L_1}
    \widetilde{\mathcal M}^{\mathrm{reg}}(p,q),
\]
其中对每个给定的 $p,q\in L_0\cap L_1$，定义
\[
  \widetilde{\mathcal M}^{\mathrm{reg}}(p,q)
  =\left\{
    u\in\widetilde{\mathcal M}^{\mathrm{reg}}(L_1,L_0)
    \ \middle|\
    \begin{array}{l}
      \displaystyle\lim_{\tau\to\infty}u(-\tau,\mathord\cdot)=\ell_p,\\[-1pt]
      \displaystyle\lim_{\tau\to\infty}u(\tau,\mathord\cdot)=\ell_q
    \end{array}
  \right\}.
\]
这里 $\widetilde{\mathcal M}^{\mathrm{reg}}(p,q)$ 扮演连接轨道空间\liftIndex{I040}的角色；也就是说，
它对应于有限维莫尔斯理论中连接光滑函数两个临界点的流线。

由于 $\{J_t\}_t$ 不依赖于 $\tau$，定义域上的 $\tau$-平移给出一个自然的
$\mathbb R$-作用。这个作用保持
$\widetilde{\mathcal M}^{\mathrm{reg}}(p,q)$，因而诱导其上的作用。对每个给定的
$p,q\in L_0\cap L_1$，把这个作用的商记为
\[
  \mathcal M^{\mathrm{reg}}(p,q)
  =\widetilde{\mathcal M}^{\mathrm{reg}}(p,q)/\mathbb R.
\]
事实上，空间 $\mathcal M^{\mathrm{reg}}(p,q)$ 还可以按照其元素的同伦类进一步分解。
为了恰当地编码冒泡现象，研究这个分解非常重要。

\begin{definition}{2.3.7}\label{def:2-3-7}
设 $p,q\in L_0\cap L_1$。以
$\pi_2(p,q)=\pi_2(p,q;L_0,L_1)$ 表示满足边界条件
\[
  u(0,t)\equiv p,
  \qquad u(1,t)=q;
  \qquad u(s,0)\in L_0,
  \qquad u(s,1)\in L_1
\]
的光滑映射 $u:[0,1]\times[0,1]\to M$ 相对于边界的同伦类集合。以
$[u]\in\pi_2(p,q)$ 表示 $u$ 的同伦类，并以 $B$ 表示
$\pi_2(p,q)$ 的一般元素。

又以 $\pi_2(p)$ 表示对子 $(\ell_p,w)$ 的同伦类集合；也就是说，它是满足
$w(0,\mathord\cdot)=\ell_0$、$w(1,\mathord\cdot)=\ell_p$，且对所有
$s\in[0,1]$ 都有 $w(s,\mathord\cdot)\in\Omega(L_0,L_1)$ 的映射
$w:[0,1]^2\to M$ 相对于两端的同伦类集合。

对给定的 $B\in\pi_2(p,q)$，以 $\operatorname{Map}(p,q;B)$ 表示类 $B$
中这种 $u$ 所成的集合。每个元素 $B\in\pi_2(p,q)$ 都由显然的粘合映射
$w\mapsto w\#u$ 诱导一个映射
\[
  (\mathord\cdot)\#B:\pi_2(p)\longrightarrow\pi_2(q),
\]
其中 $u\in\operatorname{Map}(p,q;B)$。具体地说，我们固定用下式定义
$w\#u:[0,1]^2\to M$：
\[
  (w\#u)(s,t)=
  \begin{cases}
    w(2s,t),&0\leq s\leq\dfrac12,\\[4pt]
    u(2s-1,t),&\dfrac12\leq s\leq1.
  \end{cases}
\]
此外还有自然的粘合映射
\[
  \pi_2(p,q)\times\pi_2(q,r)\longrightarrow\pi_2(p,r),
\]
它由拼接 $(u_1,u_2)\mapsto u_1\#u_2$ 诱导。映射
$u_1\#u_2:[0,1]^2\to M$ 也一律用上面的方式明确表示。
\end{definition}

由 $L_0,L_1$ 的拉格朗日性质，下面的引理是标准的。

\begin{lemma}{2.3.8}\label{lem:2-3-8}
设 $p,q\in L_0\cap L_1$。则
$u\in\operatorname{Map}(p,q;B)$ 的辛面积为常数。把
$\operatorname{Map}(p,q;B)$ 中各元素的这个公共面积记作 $\omega(B)$。
\end{lemma}

对子 $(p,q;B)$ 还相伴着一个良定义的相对指标概念，即马斯洛夫--维泰博指标
$\mu(p,q;B)$。其定义见\textbf{[Vit87,Flo88I]}。下面说明绝对指标
$\mu([\ell_p,w])$ 与相对指标 $\mu(p,q;B)$ 之间的关系。

\begin{statement}{定义兼命题}{2.3.9}\label{defprop:2-3-9}
（马斯洛夫--维泰博指标）设 $u\in\operatorname{Map}(p,q;B)$。差
\[
  \mu([\ell_p,w])-\mu([\ell_q,w\#B])
\]
只依赖于 $(p,q)$ 和同伦类 $B$，而不依赖于 $w$ 的选择。事实上，
\[
  \mu([\ell_q,w\#B])-\mu([\ell_p,w])=\mu(p,q;B).
\]
\end{statement}

定义兼命题 2.3.9 由第 3 章的命题 3.7.59（式 (3.7.60.1)）推出。
在横截情形，马斯洛夫--维泰博指标控制弗洛尔模空间的维数。

\begin{definition}{2.3.10}\label{def:2-3-10}
记
\[
  \begin{split}
  \widetilde{\mathcal M}^{\mathrm{reg}}(p,q;B)
  &=\widetilde{\mathcal M}^{\mathrm{reg}}
      (p,q;B;\{J_t\}_t)\\
  &=\widetilde{\mathcal M}^{\mathrm{reg}}(L_1,L_0,p,q;B)
   =\widetilde{\mathcal M}^{\mathrm{reg}}
      (L_1,L_0,p,q;B;\{J_t\}_t)
  \end{split}
\]
为柯西--黎曼方程~\eqref{eq:2-3-3} 满足渐近条件与同伦条件
\[
  u(-\infty,\mathord\cdot)\equiv\ell_p,
  \qquad
  u(\infty,\mathord\cdot)\equiv\ell_q,
  \qquad
  [u]=B
\]
的有限能量解集合。

这里要说明，虽然 $u$ 先验地定义在 $\mathbb R\times[0,1]$ 上，但当
$L_0,L_1$ 横截相交时，由柯西--黎曼方程之解的指数衰减性质，它可以紧化成具有相应边界条件的连续映射
$\overline u:[0,1]\times[0,1]\to M$。例如，可以把 $\overline u$ 定义为复合
\[
  \overline u(s,t)=u(h^{-1}(s),t),
\]
其中 $h:\mathbb R\to(0,1)$ 是由下式定义的函数
\translatornote{原书把一元函数 $h(\tau)$ 的类型写成积空间之间的映射；紧前复合与紧后逆函数都要求 $h:\mathbb R\to(0,1)$。}：
\[
  h(\tau)=\frac12\left(\frac{\tau}{\sqrt{1+\tau^2}}+1\right).
\]
$u$ 的指数衰减性质蕴含 $\overline u$ 连续延拓到 $[0,1]\times[0,1]$。
我们称 $\overline u$ 为 $u$ 的紧化映射。沿用记号，也以 $[u]$ 表示紧化映射
$\overline u$ 的类 $[\overline u]\in\pi_2(p,q)$。
\end{definition}

下面这个基本的横截性结论与维数结论由弗洛尔在\textbf{[Flo88I,88II]} 中证明。
关于完成弗洛尔横截性证明时所需的弗洛尔流线一般结构的若干技术性结果，另见
\textbf{[FHS95,Oh93,97I]}。

\begin{theorem}{2.3.11}\label{thm:2-3-11}
（弗洛尔，\textbf{[Flo88I,88II]}）设
$J_0,J_1\in\mathcal J_\omega$ 是任意一对固定的相容殆复结构。则存在稠密子集
\[
  \mathcal P(\mathcal J_\omega)^{\mathrm{reg}}(L_0,L_1;J_0,J_1)
  \subset\mathcal P(\mathcal J_\omega;J_0,J_1),
\]
使得对每个
$\{J_t\}_t\in
\mathcal P(\mathcal J_\omega)^{\mathrm{reg}}(L_0,L_1;J_0,J_1)$，
以及任意 $p,q\in L_0\cap L_1$ 和 $B\in\pi_2(p,q)$，模空间
\[
  \widetilde{\mathcal M}^{\mathrm{reg}}
    (L_1,L_0,p,q;B;\{J_t\}_t)
\]
在非空时横截，且为维数 $\mu(p,q;B)$ 的光滑流形。
\end{theorem}

\subsection{弗洛尔的定义：\texorpdfstring{$\mathbb Z_2$}{Z2}-系数}
\label{sec:2-3-2}
\liftIndex{I049}

本小节回顾弗洛尔对对子 $(L_0,L_1)$ 取 $\mathbb Z_2$-系数的弗洛尔上同调的原始定义
\textbf{[Flo88IV]}。整数系数的定义涉及非平凡的定向问题；对它的研究将推迟到第 8 章。

\begin{definition}{2.3.12}\label{def:2-3-12}
设 $L_0,L_1$ 横截。记
\[
  \operatorname{Int}(L_1,L_0;\ell_0)
  =\{p\in L_0\cap L_1
    \mid \ell_p\text{ 位于 }\Omega(L_0,L_1;\ell_0)\},
\]
并记
\[
  CF_{\mathbb Z_2}(L_1,L_0;\ell_0)
  =\text{以 }\operatorname{Int}(L_1,L_0;\ell_0)
    \text{ 为基的 }\mathbb Z_2\text{-向量空间}.
\]
假设 $\{J_t\}_t\in\mathcal P(\mathcal J_\omega)$ 满足以下条件。

\par\smallskip\noindent\textup{(2.3.12.1)}\quad
对满足 $\mu(p,q;B)\leq0$ 的任意
$p,q\in\operatorname{Int}(L_1,L_0;\ell_0)$，除非
$p=q$ 且 $B=0\in\pi_2(p,q)$，否则
\[
  \widetilde{\mathcal M}(p,q;B;\{J_t\}_t)=\varnothing.
\]
当 $p=q$ 且 $B=0$ 时，唯一的解是定常解；也就是说，对所有
$(\tau,t)\in\mathbb R\times[0,1]$，都有 $u(\tau,t)\equiv p=q$。

\par\smallskip\noindent\textup{(2.3.12.2)}\quad
对任意 $p,q\in\operatorname{Int}(L_1,L_0;\ell_0)$ 以及满足
$\mu(p,q;B)=1$ 的同伦类 $B\in\pi_2(p,q)$，模空间
\[
  \widetilde{\mathcal M}(p,q;B;\{J_t\}_t)/\mathbb R
\]
是光滑流形（弗雷德霍姆正则）。此外，商空间
\[
  \mathcal M(p,q;B;\{J_t\}_t)
  =\widetilde{\mathcal M}(p,q;B;\{J_t\}_t)/\mathbb R
\]
是紧的，因而是有限集。记
\[
  n(p,q;B)=\#\mathcal M(p,q;B;\{J_t\}_t)\pmod2
\]
为 $\mathcal M(p,q;B;\{J_t\}_t)$ 中元素的模 $2$ 计数。当
$\mu(p,q;B)\neq1$ 时，令
\[
  n(p,q;B)=0.
\]

\par\smallskip\noindent\textup{(2.3.12.3)}\quad
对任意 $p,r\in\operatorname{Int}(L_1,L_0;\ell_0)$ 以及满足
\[
  \mu(p,r;B)=2
\]
的 $B\in\pi_2(p,r)$，空间 $\mathcal M(p,r;B;\{J_t\}_t)$ 可紧化为一个光滑带边
一维流形，其边界只由断裂流线 $[u_1]\#_\infty[u_2]$ 组成；这里对所有可能的
$q\in\operatorname{Int}(L_1,L_0;\ell_0)$ 以及满足下列条件的
$B_1\in\pi_2(p,q)$、$B_2\in\pi_2(q,r)$，有
$u_1\in\widetilde{\mathcal M}(p,q;B_1;\{J_t\}_t)$、
$u_2\in\widetilde{\mathcal M}(q,r;B_2;\{J_t\}_t)$，且
\[
  B_1\#B_2=B;
  \qquad
  u_1\in\widetilde{\mathcal M}(p,q;B_1;\{J_t\}_t),
  \qquad
  u_2\in\widetilde{\mathcal M}(q,r;B_2;\{J_t\}_t),
\]
以及
\[
  \mu(p,q;B_1)=\mu(q,r;B_2)=1.
\]
称任何这样的 $\{J_t\}_t$ 是 $(L_0,L_1)$-正则的，并称任何这样的三元组
$(L_0,L_1;\{J_t\}_t)$ 是弗洛尔正则的。
\translatornote{原书在本定义中先后出现四类互相冲突的记号：链群的基把 $L_0,L_1$ 倒序；点对误写成子集；(2.3.12.2) 的双重否定会在指标 $1$ 时把计数置零；(2.3.12.3) 又让外端点 $q$ 同时充当断裂中间点。译文依本节既定次序和断裂流线 $p\to q\to r$ 的类型分别补正。}
\end{definition}

这个定义中的主要假设，涉及 $\mathcal M(L_1,L_0;\{J_t\}_t)$ 的零维分支
（对应 (2.3.12.2)）和一维分支（对应 (2.3.12.3)）的某些横截性与紧性性质。

弗洛尔在\textbf{[Flo88II,88IV]} 中证明：若
$L_0=L$、$L_1=\psi_1^H(L)$，且拉格朗日子流形对子满足
$\pi_2(M,L)=\{e\}$，则这些假设对 $\mathcal P(\mathcal J_\omega)$ 的某个稠密子集中的
相容殆复结构族 $\{J_t\}_t$ 成立；这里 $\psi_1^H$ 是由时间依赖函数
$H:[0,1]\times M\to\mathbb R$ 生成的哈密顿微分同胚。在这种情形，假设 $L$ 连通，
弗洛尔把讨论限制在 $\Omega(L_0,L_1)$ 中包含道路 $\ell_{\mathrm{can}}$ 的典范连通分支；该道路定义为
\[
  \ell_{\mathrm{can}}(t)=\psi_t^H(p_0),
\]
其中 $p_0\in L$ 是一个固定点。

不难验证，当 $\pi_2(M,L)=0$ 时，对每个给定的
$p,q\in L_0\cap L_1$，集合 $\pi_2(p,q)$ 都含有唯一元素，因此在以上所定义各种对象的记号中
省略对 $B$ 的依赖。特别地，作用量泛函 $\mathcal A$ 和马斯洛夫--莫尔斯指标 $\mu$
已经在道路空间 $\Omega(L_0,L_1;\ell_{\mathrm{can}})$ 上良定义（至多相差一个普适的常数平移），
因而表述弗洛尔上同调时不需要诺维科夫环。所以在这种情形，这一定义与有限维莫尔斯理论
\textbf{[Mil63,Flo89II]} 十分相似。

以上定义的要点是：对弗洛尔正则三元组
$(L,\psi_1^H(L);\{J_t\}_t)$，弗洛尔上边界映射
\[
  \delta_0:CF_{\mathbb Z_2}(\psi_1^H(L),L;\ell_{\mathrm{can}})
  \longrightarrow
  CF_{\mathbb Z_2}(\psi_1^H(L),L;\ell_{\mathrm{can}})
\]
由矩阵元 $n(p,q)$ 定义
\translatornote{原书在此重新以输入 $p$ 为求和变量，并留下游离的 $q$；译文依从 $p$ 到 $q$ 的矩阵元约定改为对 $q$ 求和并输出基元 $q$。}，即
\[
  \delta_0p
  =\sum_{q\in L\cap\psi_1^H(L)}n(p,q)q,
\]
并满足 $\delta_0\delta_0=0$。

于是，弗洛尔正则三元组
$(L,\psi_1^H(L);\{J_t\}_t)$ 的弗洛尔上同调定义为
\begin{equation}
  HF^*(\psi_1^H(L),L;\{J_t\}_t)
  :=\operatorname{Ker}\delta_0/\operatorname{Im}\delta_0.
  \tag{2.3.13}\label{eq:2-3-13}
\end{equation}
弗洛尔还在\textbf{[Flo88IV]} 中证明
\translatornote{原书把本节的柯西--黎曼方程错引为 (4.3)，并在时间一弗洛尔群中留下自由变量 $t$；这里分别改引 (2.3.3) 并取 $\operatorname{Graph}(\epsilon df)=\psi_1^H(o_N)$。}，改变哈密顿量 $H$ 以及殆复结构族
$\{J_t\}_t$ 时，这个上同调仍保持同构。他在\textbf{[Flo88IV]} 中关于这一不变性的原始证明使用分岔论证。
后来，\textbf{[Oh93]} 用更自然的配边论证取而代之；该论证使用带移动拉格朗日边界条件
\[
  u(\tau,0)\in\psi_{\chi(\tau)}^H(L),
  \qquad
  u(\tau,1)\in L
\]
的柯西--黎曼方程，其中 $\chi:\mathbb R\to[0,1]$ 是光滑函数，并且当 $\tau$
充分小时 $\chi(\tau)=0$，当 $\tau$ 充分大时 $\chi(\tau)=1$。而这一做法又受到弗洛尔在
\textbf{[Flo89I]} 中针对周期轨道情形的类似构造启发。

随后，弗洛尔在\textbf{[Flo89II]} 中研究了余切丛
$\pi:T^*N\to N$ 上零截面 $L=o_N$ 的情形，所取对子
$(\{J_t\}_t,H)$ 为
\[
  H=(\epsilon f)\circ\pi,
  \qquad
  J_t\equiv(\psi_t^H)_*J_g,
\]
其中 $J_g$ 是由底空间 $N$ 上给定黎曼度量 $g$ 的列维--奇维塔联络在
$T^*N$ 上诱导的典范殆复结构，而 $(f,g)$ 是一个莫尔斯--斯梅尔对。
在这种情形，
\[
  \psi_t^H(o_N)=\operatorname{Graph}(\epsilon t\,df),
\]
即 $1$-形式 $\epsilon t\,df$ 的图像。

对充分小的 $\epsilon>0$，弗洛尔建立了式~\eqref{eq:2-3-3} 的解空间
$\mathcal M(p,q)$ 与满足梯度流线方程
\[
  \dot\chi-\operatorname{grad}_g f(\chi)=0
\]
且始于 $p$、终于 $q$ 的流线空间之间的一一对应。再利用
\textbf{[Mil63,Wit82,Con78]} 中的莫尔斯同调论证，他证明
\[
  HF^*(\operatorname{Graph}(\epsilon\,df),o_N;\{J_t\}_t)
  \cong H^*(N;\mathbb Z_2).
\]

弗洛尔继而断言，通过达布--魏因斯坦定理，上述余切丛上的考察蕴含：在
$\pi_2(M,L)=\{e\}$ 的情形，
$HF^*(\psi_1^H(L),L;\{J_t\}_t)$ 同构于 $H^*(L;\mathbb Z_2)$。
不过，我们想指出，对满足 $\pi_2(M,L)=\{e\}$ 的一般 $(M,L)$，证明的最后一步并非完全平凡。
这是因为需要证明：当 $\psi=\psi_1^H$ 收敛到恒等映射时，
$\mathcal M(\psi(L),L;\{J_t\}_t)$ 的所有元素都变薄，从而它们的像都包含在一个
达布--魏因斯坦邻域中。在弗洛尔所考察的 $\pi_2(M,L)=\{e\}$ 情形，把弗洛尔流线空间的少量收敛论证
与一致 $C^1$-界结合起来，确实不难证明这一点；后一界源于假设
$\pi_2(M,L)=\{e\}$ 排除了冒泡现象。关于这一点的详细说明见\textbf{[Oh96I]}。
下一节引入弗洛尔上同调的博特--莫尔斯版本时，我们将回到这些问题，并说明它与
$\psi\neq\operatorname{id}$ 时的 $HF^*(\psi(L),L;\{J_t\}_t)$ 之间的关系。

\subsection{博特--莫尔斯弗洛尔上同调}
\label{sec:2-3-3}
\liftIndexRangeStart{I042}
\liftIndex{I048}

考虑哈密顿同痕的对子 $(L_0,L_1)$，其中
\[
  L_0=L,
  \qquad
  L_1=\psi_1^H(L),
\]
而 $H:[0,1]\times M\to\mathbb R$ 是一个哈密顿量。人们自然希望利用
$HF(\psi_1^H(L),\allowbreak L;\{J_t\}_t)$ 的不变性，并像弗洛尔处理余切丛零截面那样，通过选取特殊类型的对子
$(\{J_t\}_t,H)$ 来计算弗洛尔上同调
$HF(\psi_1^H(L),L;\{J_t\}_t)$。然而，当图景中存在非常值泡时，还需要分析这些泡如何与弗洛尔边界映射相互作用。
当 $\psi^H$ 远离恒等映射，更确切地说，当 $\psi_1^H(L)$ 远离 $L$ 时，冒泡会与弗洛尔流线混合，
很难同后者解耦。另一方面，当 $\psi^H$ 在 $C^1$ 意义下接近恒等映射时，弗洛尔边界映射会发生厚薄分解：
在这个分解中，当在达布邻域内写成
\translatornote{原书把微分同胚 $\psi_1^H$ 与子流形的图像直接等同；译文补出其作用对象 $L$。}
\[
  \psi_1^H(L)=\operatorname{Graph}(\epsilon\,df),
  \qquad
  H=\epsilon f\circ\pi
\]
时，薄流线——更准确地说，像包含在 $L$ 的给定达布--魏因斯坦邻域中的流线——对应于莫尔斯梯度流线。
所有其他流线称为厚流线，它们是泡与莫尔斯梯度流线的组合。事实上，先前已有许多关于紧拉格朗日嵌入之辛拓扑的非平凡结果，
是通过研究弗洛尔上同调中的这种厚薄分解得到的。这类应用见\textbf{[Chek98,Oh96I]}；
当 $\epsilon\to0$ 时极限图景的描述见\textbf{[Fuk97III,Oh96I,96II]}。
对这种绝热退化作精确研究，结果是高度非平凡且技术性很强的。

人们甚至可能希望不取上述绝热极限，而直接研究退化情形
\begin{equation}
  H=0\quad\text{（或 $\psi^H=\operatorname{id}$）},
  \qquad
  J_t\equiv J_0\quad\text{与时间无关}.
  \tag{2.3.14}\label{eq:2-3-14}
\end{equation}
事实上，弗洛尔上同调的定义可以推广到干净相交的对子 $(L_0,L_1)$。
在这种情形，我们的泛函 $\mathcal A$ 是博特--莫尔斯泛函\textbf{[Bot59]} 的一个无限维类比。
例如，可以直接考虑对子 $L_0=L_1=L$。这一构造的细节见第 3.7.5 小节。

弗洛尔同调中的博特--莫尔斯理论，最早由第一署名作者深谷贤治于 1992 年在规范场论背景下
\textbf{[Fuk96II]} 完成。（在\textbf{[Fuk96II]} 之前发表的托马什·姆罗夫卡（Tomasz Mrowka）的工作
\textbf{[Mro89]} 也与此密切相关。）在哈密顿微分同胚的弗洛尔理论背景下，
谢尔盖·亚历山德罗维奇·皮乌尼欣（Sergey Aleksandrovich Piunikhin）
在\textbf{[Piu94]} 中提出了这一构造，阮勇斌（Yongbin Ruan）--田刚（Gang Tian）也在
\textbf{[RuTi95]} 中提出了它；皮乌尼欣--萨拉蒙--马蒂亚斯·施瓦茨（Matthias Schwarz）
则在\textbf{[PSS96]} 中概述了一种略有不同的进路，其中使用的模空间类似于
\textbf{[Fuk97III]} 中引入的模空间。（据\textbf{[PSS96]} 第 171 页第 11 行，
\textbf{[PSS96]} 并未给出分析细节。）本书第 7 章第 7.1、7.2 节提供了补全
\textbf{[Piu94]} 与\textbf{[RuTi95]} 所述博特--莫尔斯理论之分析细节所需的材料。

近期，Ke Zhu 与第二署名作者吴龙根在\textbf{[OhZhu07]} 中发展了尺度依赖粘合；
它给出了\textbf{[PSS96]} 所采用进路的分析细节。尽管如此，也可以不重新缩放目标流形，
而在博特--莫尔斯背景下沿用弗洛尔粘合的更标准进路写出证明。在没有冒泡现象时，
M. Pozniak 先前已经针对拉格朗日相交开展过一项相关研究\textbf{[Poz99]}
（更多讨论见第 7.2.2 小节）。

直接在博特--莫尔斯情形~\eqref{eq:2-3-14} 下工作的一个优点是，可以把冒泡现象或量子贡献同上同调的经典部分完全解耦。
这里我们使用奇异上同调而不是莫尔斯上同调；若对非零莫尔斯函数 $f$ 采用
$\epsilon\to0$ 的绝热退化进路，则后者会自然出现。这样，在一些有利情形中，
就可以把弗洛尔上同调的计算约化为对全纯圆盘的研究。Cheol-Hyun Cho 与第二署名作者曾利用这一思想，
计算法诺环面流形的拉格朗日环面纤维之弗洛尔上同调\textbf{[Cho04I,ChOh03]}；
这在数学上证明了堀健太朗于\textbf{[Hor01]} 中作出的镜像对称预言，该预言涉及环面流形上的
$A$-模型及其朗道--金兹堡 $B$-模型。

我们在第 3 章构造的、与拉格朗日子流形相伴的 $A_\infty$ 代数，可以看作弗洛尔理论的博特--莫尔斯版本
（带乘积结构）；在其中，拉格朗日子流形的所有对象同时塌缩为单个拉格朗日子流形。特别地，
只要其 $m_1$-上同调有定义，它就同构于弗洛尔上同调
$HF^*(\psi_1^H(L),L)$。后一陈述的证明见第 5 章第 5.3 节。\liftIndexRangeEnd{I042}
